\documentclass[../main.tex]{subfiles}
\graphicspath{{\subfix{../images/}}}

\begin{document}
\subsection{Programació no lineal}
\chapter{Teoria d'extrems en R elevat a n}
\chapter{El problema general d'optimització en R elevat a n}
\begin{displaymath}
\min_{x \in S} f(x)
\end{displaymath}
$S \subset \mathbb{R}^n$, $f$ a valors en $\mathbb{R}$.
\begin{definicio}
	$x^* \in S$ és un mínim global (o absolut) de $f$ en $S$ si $\forall x \in S,\;f(x)\geq f(x^*)$
\end{definicio}

\begin{definicio}
	$x^* \in S$ és un mínim local de $f$ en $S$ si $\exists \varepsilon > 0$ tal que $\forall x \in S \cap B(x^*, \varepsilon),\;f(x) \geq f(x^*)$
\end{definicio}

\begin{obs}
\begin{displaymath}
\max f = -\min(-f)
\end{displaymath}
\end{obs}
\begin{teorema}
	Sí $f$ es continua i $S$ es compacte a $\mathbb{R}^n$ llavors $\exists$ mínim global de $f$ en $S$
\end{teorema}

\begin{definicio}
\begin{itemize}
\item $S$ és el conjunt factible
\item $x \in S$ són els punts factibles
\item $x \not\in S$ són els punts no factibles
\end{itemize}
\end{definicio}
\begin{notacio}
Vector gradient en $x \in \mathbb{R}^n$ és $\nabla f(x) = \left(\frac{\partial f}{\partial x_1}(x), \dots, \frac{\partial f}{\partial x_n}(x)\right)$\\
Derivada direccional és $\nabla f(x) \cdot d = \lim\limits_{h\to0}\frac{f(x+hd)-f(x)}{h}$\\
Matriu hessiana en $x \in \mathbb{R}^n$ és $\nabla^2 f(x) = \begin{pmatrix}\frac{\partial^2 f}{\partial^2 x_1}(x) & \dots & \frac{\partial^2 f}{\partial x_1 \partial x_n}(x)\\\vdots & \ddots & \vdots\\ \frac{\partial^2 f}{\partial x_n \partial x_1}(x) & \dots & \frac{\partial^2 f}{\partial^2 x_n}(x)\end{pmatrix}$
\end{notacio}
\begin{teorema}[Condicions necessaries de mínim local]
$x^*$ és el mínim local de $f$ en $S$, aleshores $\forall d$ factible en $x^*$
\begin{itemize}
\item $\nabla f(x^*) \cdot d \geq 0$
\item Si $\nabla f(x^*) \cdot d = 0 \Rightarrow d^t\cdot\nabla^2f(x^*)\cdot d \geq 0$
\end{itemize}
\end{teorema}
Si $x^*$ és interior de $S$,
\begin{itemize}
\item $\nabla f(x^*) = 0$
\item $\nabla^2 f(x^*)$ és semidefinida positiva
\end{itemize}
\begin{definicio}
Quan $d^tMd \geq 0\;\forall d \in \mathbb{R}^n\setminus\{0\}$ diem que $M$ és semidefinida positiva.
\end{definicio}
\begin{obs}
Taylor (les sèries, no la cantant) en una dimensió: $f(x + h) = \sum\limits_{n\geq 0} f^{(n)}(x)\frac{h^n}{n!}$.\\
En $\mathbb{R}^n$ serà $f(x+hd) = \sum\limits_{n\geq 0} {(hd)^n}^t\nabla^{(n)}f(x)\frac{(hd)^n}{n!}$
\end{obs}

\begin{teorema}[Condicions suficients de mínim local]
	Si $x^* \in S$, i
	\begin{itemize}
		\item $\nabla f(x^*) = 0$
		\item $\nabla^2 f(x^*)$ és definida positiva
	\end{itemize}
\end{teorema}

\chapter{Convexitat}
\begin{definicio}
	Un conjunt $S$ és convex si $\forall x_1, x_2 \in S,\;\forall \lambda \in (0, 1),\;\lambda x_1 + (1-\lambda)x_2 \in S$
\end{definicio}
\begin{definicio}
	Sigui $S$ convex, $f: S \to \mathbb{R}$ és convex si $\forall x_1, x_2 \in S,\;\forall\lambda\in(0,1),\;f(\lambda x_1+(1-\lambda)x_2) \leq \lambda f(x_1)+(1-\lambda)f(x_2)$
\end{definicio}
\begin{proposicio}
$f$ és convexa en $S \iff f(x) \geq f(x_0)\nabla f(x)(x-x_0)\forall x_0,x \in S$
\end{proposicio}
\begin{proposicio}
$f$ és convexa en $S \iff \forall x \in S,\;\forall d$ factible en $x,\;d^t\nabla^2 f(x)d \geq 0$
\end{proposicio}
\begin{teorema}[Local-global per a funcions convexes]
$S\subset \mathbb{R}^n$ convex i $f : S \to \mathbb{R}$ convexa. Si $x^*$ és mínim local de $f$ en $S$, llavors és mínim global.
\end{teorema}
\begin{teorema}[Les condicions necessaires són suficients]
Si $x^*$ compleix les condicions necessàries, aleshores $x^*$ és mínim local (i global) de $f$ en $S$.
\end{teorema}
\chapter{Restriccions funcionals}
$S = $ punts que compleixen un sistema d'equacions (o inequacions) de $m$ igualtats (o desigualtats).
\begin{notacio}
	Quan $f$ i totes aquestes funcions $h$ i $g$ són afins, diem que $\min\limits_{x \in S} f(x)$ és un problema de programació lineal. Si no, tenim un problema de programació no lineal.\\
	La funció $f$ s'anomena funció objectiu i les igualtats i desigualtats són les restriccions del problema.
\end{notacio}
\begin{definicio}
	Donat $x^*$ factible, una restricció $g \leq 0$ és activa en $x^*$ si $g(x^*) = 0$. És inactiva si $g(x) < 0$
\end{definicio}
\begin{definicio}
Un vector $v$ és tangent a $S \subset \mathbb{R}^n$ en un punt $x^* \in S$ si $\exists \left\{x_k\right\} \subset S$ tal que $\lim\limits_{k\to\infty}x_k = x^*$, i una successió $\{t_k\}_k\subset\mathbb{R}^+$ decreixent cap a $0$ tals que
\begin{displaymath}
\lim\limits_{k\to x_0}\frac{x_k-x^*}{t_k} = v
\end{displaymath}
\end{definicio}
\subsubsection{Algoritmes d'optimització en R}
\chapter{Mètode de Newton}

Tenim $f: \mathbb{R} \to \mathbb{R}$ i un punt inicial $x_1 \in \mathbb{R}$. Com a funció aproximadora agafem $P(x) = f(x_1) + f'(x_1)(x-x_1)+\frac{1}{2}f''(x_1)(x-x_1)^2$\\
Fem la derivada i igualem a $0$, llavors tenim $x_2 = \frac{f'(x_1)}{f''(x_1)}$. Si això ho fem succeir, si té un límit a un punt $x_k$, aquest és un extrem local.

\begin{proposicio}
	Suposem que tenim $f$ 3 vegades derivable i $x^*$ tal que $f'(x^*) = 0$ i que $f''(x^*) \neq 0$, aleshores $\exists \delta > 0:\;|x_1-x^*|<\delta$, la successió $\left\{x_k = x_{k-1} - \frac{f'(x_k)}{f''(x_k)}\right\}$
\end{proposicio}

\chapter{Ordres de convergència}

\begin{displaymath}
|x_k|
\end{displaymath}

\end{document}
